\section{Feature-Cut Priorities}
\label{sec:feature-cuts}

This program is not optimized around producing the maximum number of PCBs. It is optimized around reaching November~30 with two validated customer architectures built on one coherent Alfred hardware platform. If schedule slips, cuts happen in this order.

\subsection{Protect at all costs}

\begin{enumerate}[itemsep=0.1em]
\item LAN8650/LAN8651 (and, time permitting, LAN8660) learning (Section~\ref{sec:lan866x}).
\item Rev~A.
\item Rev~B.
\item The common-node architecture itself (Section~\ref{sec:common-platform}).
\item Multidrop 10BASE-T1S.
\item Greenfield E2B demo (Section~\ref{sec:arch-a-demo}).
\item CAN brownfield Gateway (Section~\ref{sec:arch-b}).
\item Authorized bidirectional operation (Mode~2, Section~\ref{sec:gateway-modes}).
\item Quantitative validation (Section~\ref{sec:validation-plan}).
\end{enumerate}

\subsection{Cut first}

\begin{itemize}[itemsep=0.1em]
\item 100BASE-T1 / 1000BASE-T1.
\item Any protocol beyond CAN/CAN-FD/(LIN).
\item Sophisticated enclosures.
\item Large peripheral counts on the E2B demo (one sensor, one actuator is sufficient; more is not required).
\item Additional node SKUs beyond E2B and Gateway.
\item Cosmetic features (labeling polish, connector aesthetics) beyond what Section~\ref{sec:rev-c} already treats as minimum-credible.
\end{itemize}

\begin{decision}[title={LIN and Rev C are explicitly conditional}]
LIN demonstration becomes optional the moment it threatens CAN/CAN-FD/T1S validation (Section~\ref{sec:arch-b}) -- it is never protected at the expense of the mandatory bus set. Rev~C is dropped or simplified (fall back to demonstrating on characterized Rev~B hardware, cleaned up as much as time allows) rather than letting a Rev~C schedule slip jeopardize Rev~B validation, which is the load-bearing proof for the entire program (Section~\ref{sec:rev-c}).
\end{decision}
